\paragraph{}
Au cours de ce projet, certaines difficultés rencontrées auront parfois ralenti l'avancement initialement prévu.

\subsubsection{Implémentation des algorithmes}
\paragraph{}
L'implémentation des algorithmes d'anaglyphes et d'autostéréogrammes était le coeur de notre projet et du logiciel que nous avions à mettre en place, il était donc important que nous lui accordions un temps conséquent durant notre projet. C'est pour cela que nous avions initialement prévu de faire travailler deux personnes sur ces algorithmes durant un mois, en partant des recherches qui avaient été effectuées pour le cahier des charges. Toutefois, l'implémentation de deux algorithmes étant nécessaire pour chacun de ces rendus, ce mois a finalement été trop court pour nos équipes de programmeurs.

\paragraph{}
Pour les autostéréogrammes, le rendu donné par le premier algorithme choisi était découpé en tranches comme nous pouvoir le voir sur la figure \ref{fig:sphAutoste}.

\begin{figure}[h]               
	\centering      
	\includegraphics[scale=0.7]{bouleAutoste.png}
	\caption{\label{fig:sphAutoste} Autostéréogramme d'une boule avec un effet en tranche \protect}
\end{figure}

Cet effet avait en effet été souligné dans l'article de Witten, Inglis et Thimbleby, mais l'utilisation de points aléatoires pour le fond de l'image accentuait cet effet de tranches. 
Il a alors fallu mettre en place le second algorithme que nous avons présenté précédemment dans ce rapport, ce qui a allongé la période de réalisation. De plus, l'optimisation proposé pour ce deuxième algorithme par W.A.Steer posait dans un premier temps problème pour la génération de l'autostéréogramme. Il aura alors fallu procéder étape par étape pour obtenir dans un premier temps un rendu sans l'optimisation de l'algorithme, qui aura été ajoutée par la suite.

\paragraph{}
Pour les anaglyphes, la première difficulté était la récupération de paires d'images sur lesquelles travailler. Notre équipe a tout d'abord essayé de travailler sur la scène avec OpenGL pour récupérer des images qu'ils pourraient utiliser. Toutefois, une autre difficulté, liée aux shaders comme nous le présenterons dans la partie qui suit, les a empêché de poursuivre dans cette voie. Ils ont donc choisi de travailler sur des images trouvées sur Internet afin d'avancer dans l'implémentation. 
Une autre difficulté de l'algorithme des anaglyphes est le travail sur la saturation gamma. Pour permettre un rendu agréable à visualiser, notre équipe a cherché à travailler sur cette saturation, mais l'effet obtenu n'était pas toujours celui qu'ils espéraient. En effet, selon les couleurs supportées par les écrans d'ordinateur, il arrive que le rendu obtenu soit ``brûlé'', c'est-à-dire que comme les valeurs sont trop hautes l'ordinateur ne sait pas rendre la couleur et transforme les zones en blanc. Ce problème n'a malheureusement pas pu être résolu, car il faudrait pouvoir adapter l'algorithme aux caractéristiques de l'écran sur lequel il est testé.
Enfin, nos clients souhaitant pouvoir imprimer les anaglyphes, il était très important que les couleurs lors de l'impression révèlent aussi bien l'anaglyphe. En effet, si un écran d'ordinateur possède les composantes rouge, verte et bleue pour son affichage, la plupart des imprimantes utilisent le pattern cyan, magenta, jaune et noir. Pour éviter d'éventuels artefacts qui auraient pu être causés par la transformation automatique des couleurs, il a fallu traiter celles-ci pour que l'impression soit le plus parfaite possible. Une fois de plus, un problème irrésolvable au moment de l'impression est l'imprimante qui sera utilisée. En effet, selon les caractéristiques de celles-ci, le pattern de couleur utilisé peut modifier l'effet du rendu. Nous avons choisi de travailler les couleurs pour un rendu sur une imprimante utilisant le pattern Cyan-Magenta-Jaune-Noir, qui sont les plus répandues.

\subsubsection{L'utilisation de shaders}
\paragraph{}
On appelle shaders des programmes informatiques qui vont travailler directement sur la carte graphique ou le processeur graphique d'un ordinateur pour permettre un rendu meilleur et plus rapide. Ces shaders sont souvent utilisés en imagerie numérique, notamment par le logiciel Blender\footnotemark qui aura été un des logiciels-modèle pour la création du nôtre. L'utilisation de ces shaders nous paraissait importante pour permettre un rendu agréable dans notre logiciel, et pour atteindre nos objectifs de fluidité énoncés dans notre Cahier des charges.
\footnotetext{\url{http://www.blender.org/}}
\paragraph{}
Dans certains ordinateurs, et surtout dans les ordinateurs portables, il n'existe pas de carte graphique à proprement parler. Un chipset graphique en général moins puissant la remplace. Toutefois, ces chipsets peuvent empêcher l'utilisation de shaders trop récent.

\paragraph{}
Dans une première version de notre logiciel, nous avions souhaité utiliser la version 3.3 des shaders. Toutefois, celle-ci posait problèmes car elles n'étaient pas reconnue sous des machines Linux. La portabilité du logiciel étant primordiale, du moins sous les environnements Windows et Linux, nous avons dû effectuer une modification de l'ensemble du code déjà construit afin de passer à la version 2.0 des shaders qui est bien plus ancienne et qui permet la portabilité du logiciel sur quasiment toute les machines.

\subsubsection{Les outils utilisés pour la réalisation du projet}
\paragraph{}
Pour permettre la portabilité et la maintenabilité du projet, nous avons choisi de l'implémenter en C++11, en utilisant Qt5 et la bibliothèque OpenGL pour Qt, et en compilant avec l'outil CMake.

\paragraph{}
La communauté Internet sur la bibliothèque Qt et son module OpenGL est très importante et de nombreux problèmes rencontrés ont pu être résolus facilement grâce à des recherches sur certains forums. 
L'outil souvent utilisé dans la réalisation de projet Qt est QMake qui permet une compilation simplifiée avec Qt. Mais nous avons choisi de prendre CMake comme nous savions déjà l'utiliser.
CMake étant très peu utilisé pour la bibliothèque Qt, il est parfois difficile de trouver sur Internet de l'aide pour inclure des modules, par exemple le module QtXml. La plupart de la documentation relative à ce problème est très simplement géré dans QMake, mais il aura fallu beaucoup de recherches pour parvenir à trouver les bons noms de bibliothèques et de modules pour l'inclure avec CMake.

\subsubsection{Portabilité entre systèmes d'exploitation}
\paragraph{}
Tout au long de la réalisation du logiciel, il nous aura souvent fallu corriger des passages de code pour permettre la portabilité aussi bien sous Windows que sous Linux. Les conventions de codage nous auront aidé à assurer des noms de variables compatibles sous les deux environnements, mais cela n'aura parfois pas suffit. 
\paragraph{}
Un premier problème de portabilité s'est présenté avec les shaders, comme nous l'avons expliqué plus haut. Linux ne supportant pas toujours les cartes raphiques et les chipsets intégrés de la même façon que Windows, il a fallu utiliser une version plus récente pour s'assurer que le logiciel fonctionne.

\paragraph{}
Un autre cas a eu lieu lors de l'utilisation de threads pour la sauvegarde automatique. La fonction sleep() avait d'abord été utilisée, mais celle-ci ne fonctionnait que sous Linux ou avec le compilateur GCC sous Windows. Nous avons donc dû le remplacerpar l'appel de C++ sleep\_for, que l'on trouve dans le namespace this\_thread. 

\paragraph{}
La compatibilité a également été un souci pour le découpage des fichiers de chargement d'extension XML. Sous Windows, le stockage des flottants utilise un point pour séparer la partie entière de la partie décimale, aussi bien en lecture qu'en écriture. Sous Linux, la lecture utilise une virgule. Nous avons donc dû tenir compte de cette différence pour le découpage et l'utilisation de nos fichiers pour le chargement.
